\chapter{Cơ sở lý thuyết}

\section{Vi điều khiển ESP32}

ESP32 là vi điều khiển do Espressif Systems phát triển, tích hợp WiFi và
Bluetooth, được sử dụng rộng rãi trong các ứng dụng IoT. Dự án sử dụng
board ESP32-CAM tích hợp sẵn khe cắm camera và hỗ trợ PSRAM mở rộng.

\begin{table}[htbp]
\centering
\caption{Thông số phần cứng ESP32-CAM}
\label{tab:esp32-hardware}
\begin{tabular}{>{\raggedright\arraybackslash}p{0.3\linewidth} >{\raggedright\arraybackslash}p{0.6\linewidth}}
\toprule
\rowcolor{headerblue} \textcolor{white}{\textbf{Thành phần}} & \textcolor{white}{\textbf{Thông số}} \\
\midrule
\rowcolor{rowgray} Board & ESP32-CAM với PSRAM \\
Vi xử lý & Xtensa LX6 dual-core, 240 MHz \\
\rowcolor{rowgray} Camera & OV5640, JPEG, QVGA (320$\times$240) mặc định, hỗ trợ đến UXGA (1600$\times$1200) \\
Microphone & INMP441 (I2S digital), 16 kHz, mono \\
\rowcolor{rowgray} Kết nối & WiFi 802.11 b/g/n \\
Bộ nhớ & 520 KB SRAM + 4 MB PSRAM \\
\bottomrule
\end{tabular}
\end{table}

\subsection{Giao thức I2S}

I2S (Inter-IC Sound) là giao thức truyền thông nối tiếp đồng bộ được thiết kế
chuyên biệt cho dữ liệu âm thanh số. Giao thức sử dụng 3 đường tín hiệu:

\begin{itemize}
\item \textbf{SCK (Serial Clock):} Xung nhịp đồng bộ dữ liệu.
\item \textbf{WS (Word Select):} Chọn kênh trái (0) hoặc kênh phải (1).
\item \textbf{SD (Serial Data):} Đường dữ liệu nối tiếp.
\end{itemize}

Microphone INMP441 hỗ trợ I2S với độ phân giải 18-bit, đóng gói trong word
32-bit. ESP32 hoạt động ở chế độ I2S Master RX, nhận dữ liệu từ microphone.

\subsection{Giao thức TCP}

TCP (Transmission Control Protocol) đảm bảo truyền dữ liệu tin cậy, có thứ tự
giữa ESP32 và server. Dự án sử dụng TCP cho cả stream audio (port 3002) và
video (port 3001), đảm bảo không mất frame hay mẫu âm thanh.

\section{Mô hình khử nhiễu DTLN}

DTLN (Dual-signal Transformation LSTM Network) là mô hình khử nhiễu giọng nói
được đề xuất bởi Westhausen và Meyer (2020), thiết kế cho xử lý thời gian thực
trên thiết bị tài nguyên hạn chế.

\begin{table}[htbp]
\centering
\caption{Các thông số chính của mô hình DTLN}
\label{tab:dtln-params}
\begin{tabular}{>{\raggedright\arraybackslash}p{0.45\linewidth} >{\centering\arraybackslash}p{0.45\linewidth}}
\toprule
\rowcolor{headerblue} \textcolor{white}{\textbf{Thông số}} & \textcolor{white}{\textbf{Giá trị}} \\
\midrule
\rowcolor{rowgray} Tần số lấy mẫu & 16.000 Hz, mono \\
Kích thước khung (blockLen) & 512 mẫu (32 ms) \\
\rowcolor{rowgray} Bước dịch khung (block\_shift) & 128 mẫu (8 ms), overlap 75\% \\
Số bin tần số sau STFT & 257 (= blockLen / 2 + 1) \\
\rowcolor{rowgray} Số đơn vị LSTM mỗi lớp & 128 units \\
Số lớp LSTM mỗi giai đoạn & 2 lớp \\
\rowcolor{rowgray} Dropout & 0,25 \\
Số bộ lọc Encoder Conv1D & 256 \\
\rowcolor{rowgray} Hàm kích hoạt mask & Sigmoid {[}0, 1{]} \\
Hàm mất mát & Negative SNR \\
\rowcolor{rowgray} Tổng số tham số & $\approx$ 1,8 triệu \\
Kích thước file model & 3,9 MB (\texttt{DTLN\_vivos\_best.h5}) \\
\rowcolor{rowgray} Độ trễ xử lý & $\approx$ 32 ms (real-time capable) \\
\bottomrule
\end{tabular}
\end{table}

\subsection{Kiến trúc hai giai đoạn}

DTLN xử lý tín hiệu qua hai giai đoạn bổ sung lẫn nhau:

\textbf{Giai đoạn 1 -- Miền tần số (STFT Domain):}
Tín hiệu đầu vào được chia thành các khung 512 mẫu, áp dụng STFT để thu được
257 bin tần số. Phổ biên độ $|X(f,t)|$ được chuẩn hóa qua InstantLayerNormalization,
xử lý qua 2 lớp LSTM (128 units) với Dropout 0,25, rồi qua Dense + Sigmoid để tạo
mask $M_1 \in [0,1]$. Mask được nhân với phổ biên độ gốc, kết hợp lại với phổ pha
và biến đổi ngược (iFFT + Overlap-and-Add) để thu được tín hiệu $y_1(t)$.

\begin{table}[htbp]
\centering
\caption{Cấu trúc Separation Kernel -- Giai đoạn 1}
\label{tab:sep-kernel}
\begin{tabular}{>{\raggedright\arraybackslash}p{0.2\linewidth} >{\centering\arraybackslash}p{0.35\linewidth} >{\raggedright\arraybackslash}p{0.35\linewidth}}
\toprule
\rowcolor{headerblue} \textcolor{white}{\textbf{Thành phần}} & \textcolor{white}{\textbf{Thông số}} & \textcolor{white}{\textbf{Vai trò}} \\
\midrule
\rowcolor{rowgray} LSTM lớp 1 & 128 units, return\_seq=True & Đặc trưng tần số cục bộ \\
Dropout & 0,25 & Giảm overfitting \\
\rowcolor{rowgray} LSTM lớp 2 & 128 units, return\_seq=True & Phụ thuộc thời gian dài hạn \\
Dense & 257 đơn vị & Ánh xạ về không gian bin tần số \\
\rowcolor{rowgray} Sigmoid & Đầu ra {[}0, 1{]} & Tạo mask $M_1(f,t)$ \\
\bottomrule
\end{tabular}
\end{table}

\textbf{Giai đoạn 2 -- Miền đặc trưng (Feature Domain):}
Tín hiệu $y_1(t)$ được chiếu sang không gian 256 chiều qua Conv1D point-wise
(kernel=1, stride=1). Sau InstantLayerNormalization và Separation Kernel thứ hai
(cùng cấu trúc LSTM), Dense + Sigmoid tạo mask $M_2$ trong không gian đặc trưng.
Giải mã qua Conv1D point-wise (256→512) với causal padding, rồi Overlap-and-Add
tạo tín hiệu sạch $\hat{y}(t)$.

\begin{table}[htbp]
\centering
\caption{So sánh hai giai đoạn xử lý DTLN}
\label{tab:dual-signal}
\begin{tabular}{>{\raggedright\arraybackslash}p{0.2\linewidth} >{\centering\arraybackslash}p{0.33\linewidth} >{\centering\arraybackslash}p{0.33\linewidth}}
\toprule
\rowcolor{headerblue} \textcolor{white}{\textbf{Tiêu chí}} & \textcolor{white}{\textbf{STFT (GĐ 1)}} & \textcolor{white}{\textbf{Conv1D Encoder (GĐ 2)}} \\
\midrule
\rowcolor{rowgray} Loại biến đổi & Cố định (fixed basis) & Học từ dữ liệu \\
Chiều đầu ra & 257 bin tần số & 256 đặc trưng ẩn \\
\rowcolor{rowgray} Ý nghĩa vật lý & Rõ ràng, diễn giải được & Trừu tượng \\
Vai trò & Loại bỏ nhiễu tổng thể & Tinh chỉnh, phục hồi chi tiết \\
\bottomrule
\end{tabular}
\end{table}

\subsection{Hàm mất mát Negative SNR}

Mô hình được tối ưu hóa bằng hàm mất mát Negative SNR:
\begin{equation}
\text{Loss} = -10 \cdot \log_{10}\left(\frac{E[s^2]}{E[(s - \hat{s})^2 + \epsilon]}\right)
\end{equation}
trong đó $s$ là tín hiệu sạch, $\hat{s}$ là tín hiệu ước lượng, $\epsilon = 10^{-7}$.
Hàm mất mát này tối đa hóa tỷ số tín hiệu trên nhiễu (SNR) của đầu ra,
trực tiếp tối ưu chất lượng âm thanh.

\section{Mô hình XLM-RoBERTa}

XLM-RoBERTa (Cross-lingual Language Model -- Robustly Optimized BERT Approach)
là mô hình ngôn ngữ đa ngôn ngữ được phát triển bởi Facebook AI, huấn luyện
trên 2,5 TB dữ liệu CommonCrawl từ 100 ngôn ngữ.

\begin{itemize}
\item \textbf{Base model:} XLM-RoBERTa-base với 270M tham số.
\item \textbf{Fine-tuned:} \texttt{Lezh1n/xlm-roberta-ecommerce-classifier}
  cho phân loại 32 danh mục thương mại điện tử.
\item \textbf{Kích thước:} $\approx$ 1,1 GB, max sequence 512 tokens.
\item \textbf{Hiệu năng:} Accuracy 90,14\%, F1-score 90,00\%.
\end{itemize}

Ưu điểm của XLM-RoBERTa cho bài toán này là khả năng xử lý đa ngôn ngữ
(cross-lingual transfer), cho phép mô hình hiểu cả text tiếng Việt lẫn
tên sản phẩm tiếng Anh thường gặp trong thương mại điện tử.

\section{Các mô hình OpenAI}

\subsection{GPT-4o-transcribe}

GPT-4o-transcribe là mô hình Speech-to-Text mới nhất của OpenAI,
kế thừa và cải tiến từ Whisper. Mô hình hỗ trợ nhận dạng tiếng Việt
đa vùng miền với độ chính xác cao, hỗ trợ nhiều định dạng audio
(WAV, MP3, FLAC, OGG, WebM, M4A).

Trong dự án, lịch sử phát triển module STT trải qua:
PhoWhisper-tiny (local) → whisper-1 (OpenAI) → gpt-4o-transcribe,
với mỗi bước nâng cấp cải thiện đáng kể chất lượng nhận dạng tiếng Việt.

\subsection{GPT-4o-mini}

GPT-4o-mini là mô hình ngôn ngữ lớn (LLM) nhỏ gọn của OpenAI,
được sử dụng để sinh mô tả sản phẩm từ text nhận dạng giọng nói.
Mô hình có khả năng hiểu ngữ cảnh, xử lý text lộn xộn từ STT
(từ lấp, lặp lại, lỗi phát âm) và sinh văn bản mạch lạc theo
phong cách viết được chỉ định.
